\section{Distances Between Laws of Random Probabilities}\label{section:two_sample_distances}


In this paper, we consider two distance functions on $\probprobx$. The first, referred to as "Wasserstein over Wasserstein," is the natural lifting of the Wasserstein distance on $\probx$; it has been widely used in Bayesian nonparametrics \cite{nguyen2016borrowing} and machine learning \cite{ho2017multilevel, dukler2019wasserstein}. The second, termed the "Hierarchical Integral Probability Metric," is a very recent development defined in \cite{catalano2024hierarchical}. Let us start with the definition of Wasserstein over Wasserstein. 

Recall that the sample space $\X$ is Polish with a bounded diameter. Using the distance function $d_\X$ on $\X$, one can define the Wasserstein distance on $\probx$ by
$$
\W_1(P^1, P^2) = \inf_{\pi \in \Gamma(P^1, P^2)} \int_{\X \times \X} d_\X(x, y) \, \pi(dx, dy), \qquad P^1, P^2 \in \probx,
$$
where
$$
\Gamma(P^1, P^2) = \left\{ \pi \in \mathcal{P}(\X^2) : \pi(A \times \X) = P^1(A), \ \pi(\X \times A) = P^2(A) \quad \forall A \in \mathcal{B}(\X) \right\}.
$$
Note that this is the Wasserstein distance of order $1$. Since we do not consider other Wasserstein distances, we denote it simply by $\W$ instead of $\W_1$. 

In general, $\W$ is a distance function on the Wasserstein space of order $1$, defined as 
$$
\mathcal{P}_1(\X) = \left\{P \in \probx : \int_\X d_\X(x, x_0) P(dx) < \infty, \quad x_0 \in \X \right\}.
$$
Since we assume that $\X$ has a bounded diameter, all $P \in \probx$ belong to $\mathcal{P}_1(\X)$. The assumption that $\X$ is Polish and $d_{\X}$ is bounded implies that the metric space $(\probx, \W)$ is also Polish (see Remark 7.1.7 in \cite{luigi2008gradient}). Note also that $(\probx, \W)$ has a bounded diameter, as $d_{\X}$ is bounded. Therefore, instead of defining a distance function on $\mathcal{P}_1(\probx)$, we can define a distance function on $\probprobx$ using $\W$ as the ground metric: 

\begin{definition}[Wasserstein over Wasserstein]
    \label{def:wow}
    Given $Q^1, Q^2 \in \probprobx$, the Wasserstein over Wasserstein (WoW) distance between $Q^1$ and $Q^2$ is 
    $$
    \WW(Q^1, Q^2) = \inf_{\pi \in \Gamma(Q^1, Q^2)} \int_{\probx^2} \W(P^1, P^2) \, \pi(dP^1, dP^2).
    $$
\end{definition}
Here, $\Gamma(Q^1,Q^2)$ denotes the set of probability measures on $\probx \times \probx$ with marginals $Q^1$ and $Q^2$. Again, Remark 7.1.7 in \cite{luigi2008gradient} and the fact that $\W$ has a bounded diameter imply that the metric space $(\probprobx, \WW)$ is a Polish space with a bounded diameter.

The second distance function we consider is the Hierarchical Integral Probability Metric (HIPM), introduced in \cite{catalano2024hierarchical}. The primary intuition behind its definition is as follows: let $Q^1, Q^2$ be probability measures on $\probx$ and let $\widetilde{P}^1 \sim Q^1, \widetilde{P}^2 \sim Q^2$. We compute the distance between $Q^1$ and $Q^2$ by considering the distances between the random integrals $\widetilde{P}^1(f)$ and $\widetilde{P}^2(f)$, where $f$ ranges over a class of functions from $\X$ to $\R$. By considering these projections, we shift the focus from $\probprobx$ to $\mathcal{P}(\R)$, the latter of which is significantly easier to handle. As noted in \cite{catalano2024hierarchical}, such a projection is justified by the fact that for a random probability measure $\widetilde{P}$, the laws of the random integrals $\widetilde{P}(f)$, where $f$ belongs to the set of continuous and bounded functions, are sufficient to characterize the law of $\widetilde{P}$. This leads to the following definition:

\begin{definition}[Hierarchical integral probability metric]
    \label{def:hipm}
    Given $Q^1, Q^2 \in \probprobx$, the hierarchical integral probability metric (HIPM) between $Q^1$ and $Q^2$ is defined as 
    \[
    d_{\F}(Q^1, Q^2) = \sup_{f \in \F} \W(\mathcal{L}(\widetilde{P}^1(f)), \mathcal{L}(\widetilde{P}^2(f))),
    \] 
    where $\widetilde{P}^k \sim Q^k$ for $k=1,2$ and $\F$ is a class of bounded and measurable functions from $\X$ to $\R$. 
\end{definition}

Alternatively, the HIPM can be expressed as 
\[
d_{\F}(Q^1, Q^2) = \sup_{f \in \F} \W(Q^1 \circ \hat{f}^{-1}, Q^2 \circ \hat{f}^{-1}),
\]
where $\hat{f}$ denotes the lifting of the function $f$, i.e., $\hat{f}(P) = \int_\X f \, dP$, and $Q^k \circ \hat{f}^{-1}$ is the push-forward measure.

As long as the function class $\F$ is sufficiently rich, this construction defines a valid distance on $\probprobx$. Hereafter, we set $\F = \operatorname{Lip}_1(\X, \R)$ and denote the corresponding metric $d_\F$ by $\dlip$.

Both distance functions possess desirable topological properties. In particular, by Theorem 3.1 in \cite{catalano2024hierarchical}, WoW metrizes weak convergence on $\probprobx$, and HIPM does the same when $\X$ is compact.

There is a relation between the two distance functions; specifically, they are special cases of an Integral Probability Metric defined on a generic Polish space.
\begin{definition}[Integral Probability Metric]
    \label{def:ipm}
    Let $(\mathbb{Y},d_{\mathbb{Y}})$ be a Polish space and $\mathcal{H}$ a class of bounded and measurable functions from $\mathbb{Y}$ to $\R$. The Integral Probability Metric (IPM) between $P^1, P^2 \in \mathcal{P}\left(\mathbb{Y}\right)$ is 
    $$
    \mathcal{I}_{\mathcal{H}}(P^1,P^2) = \sup_{f\in\mathcal{H}} \left| \int_{\mathbb{Y}} f(y) P^1(dy) - \int_{\mathbb{Y}} f(y) P^2(dy) \right|. 
    $$
\end{definition}
Firstly, let us show that WoW can be written as an integral probability metric. By Remark 6.5 \cite{villani2008optimal}, if $P^1,P^2$ belong to $\mathcal{P}_1(\X)$, the Wasserstein space of order 1, the Wasserstein distance can be written as
$$
\W(P^1,P^2) = \sup_{f\in \text{Lip}_1(\X,\R)} \left| \int_\X f(x) P^1(dx) - \int_\X f(x) P^2(dx) \right|,
$$
where $\text{Lip}_1(\X,\R)$ denotes the set of $1-$Lipschitz continuous functions from $\X$ to $\R$. Since $(\probprobx, \WW)$ is a Polish space with a bounded diameter, we can use the above result to obtain, for all $Q^1,Q^2\in\probprobx$,
\[
\WW(Q^1,Q^2) = \sup_{f\in \text{Lip}_1(\probx,\R, \W)} \left| \int_{\probx} f(P) Q^1(dP) - \int_{\probx} f(P) Q^2(dP) \right|,
\]
where $\text{Lip}_1(\probx,\R, \W)$ denotes the set of $1-$Lipschitz continuous functions from $\probx$ to $\R$ with respect to $\W$.

Similarly, we show that HIPM can be written as an integral probability metric. For $Q^1,Q^2 \in \probprobx$,
\begin{align*}
\dlip(Q^1,Q^2) &= \sup_{f\in \text{Lip}_1(\X,\R)} \W(Q^1\circ \hat{f}^{-1}, Q^2\circ \hat{f}^{-1}) \\
&= \sup_{f\in \text{Lip}_1(\X,\R)}\sup_{g\in\lipR}\left| \int_{\R}g(x) (Q^1\circ\hat{f}^{-1})(dx) - \int_{\R}g(x) (Q^2\circ\hat{f}^{-1})(dx) \right| \\
& = \sup_{f\in \text{Lip}_1(\X,\R)} \sup_{g\in\lipR}\left|\int_{\probx}g(\widehat{f}(P))Q^1(dP) - \int_{\probx}g(\widehat{f}(P))Q^2(dP) \right| \\
& = \sup_{f\in \text{Lip}_1(\X,\R)}\sup_{g\in\lipR}\left|\int_{\probx}g(P(f))Q^1(dP) - \int_{\probx}g(P(f))Q^2(dP) \right| \\
& = \sup_{h\in \D}\left|\int_{\probx}h(P)Q^1(dP) - \int_{\probx}h(P)Q^2(dP) \right|,
\end{align*}
where $\D = \{h:\probx\to \R: h(P) = g(P(f)), \ g\in\lipR, f\in \text{Lip}_1(\X,\R) \}$.

By representing both distance functions as Integral Probability Metrics, we see that the HIPM is upper bounded by the WoW distance. Indeed, we can show that $\D$ is a subset of $\text{Lip}_1(\probx, \R, \W)$: let $h \in \mathcal{D}$ and $P^1, P^2 \in \probx$, then
\begin{align*}
|h(P^1) - h(P^2)| = |g(P^1(f)) - g(P^2(f))| \leq |P^1(f) - P^2(f) | &\leq \sup_{f \in \text{Lip}_1(\X, \R)} |P^1(f) - P^2(f) | \\
&= \W(P^1, P^2).
\end{align*}

\begin{remark}\label{rem::converg_rate}
The hierarchical estimator $Q_{n,m}$ is used to estimate the law of the random probability measure $Q$. The effectiveness of this estimator can be studied using either WoW or HIPM. A key difference between these distance functions lies in the rate of convergence of the expected distance between $Q_{n,m}$ and $Q$. In particular, by Theorem 4.1 in \cite{catalano2024hierarchical}, if $\X$ is a bounded subset of $\R^d$, then 
$$
\E \WW(Q_{n}, Q) = \mathcal{O} \left(\frac{\log \log (n)}{\log(n)}\right),
$$
where $\widetilde{P}_1, \dots, \widetilde{P}_n \overset{\mathrm{i.i.d.}}{\sim} Q$ and $Q_n = \frac{1}{n}\sum_{i=1}^n \delta_{\widetilde{P}_i}$. Note that since we only observe $Q_{n,m}$, its convergence rate will be slower than that of $Q_n$. Moreover, the theorem shows that if $Q$ is a Dirichlet process, the convergence is faster than any polynomial rate.

On the other hand, using the HIPM, we obtain a faster convergence rate. In particular, by Theorem 4.3 in \cite{catalano2024hierarchical}, if $\X$ is a bounded subset of $\R^d$, then
$$
\E \dlip (Q_{n,m}, Q) = 
\begin{cases}
\mathcal{O}\left(\frac{1}{\sqrt{n}} + \frac{1}{\sqrt{m}} \right), & \text{if } d = 1 \\
\mathcal{O}\left(\frac{1}{\sqrt{n}}\log(n) + \frac{1}{\sqrt{m}}\log(m) \right), & \text{if } d = 2 \\
\mathcal{O}\left(\frac{1}{n^{\frac{1}{d}}} + \frac{1}{m^{\frac{1}{d}}} \right), & \text{if } d > 2.
\end{cases}
$$
\end{remark}
Such rates of convergence play an important role in studying the properties of the thresholds in Section \ref{section:rademacher_complexity}.

To summarize, we consider the following hypothesis testing problem. Let $d$ denote either the WoW or the HIPM distance, and let $Q^1, Q^2 \in \probprobx$. We test the null hypothesis
\[
H_0 : d(Q^1, Q^2) = 0
\]
against the alternative
\[
H_1 : d(Q^1, Q^2) > 0.
\]
The test is based on independent samples
\[
\widehat{P}^1_{1,m}, \dots, \widehat{P}^1_{n,m} \overset{\mathrm{i.i.d.}}{\sim} Q^{1,(m)}, 
\qquad
\widehat{P}^2_{1,m}, \dots, \widehat{P}^2_{n,m} \overset{\mathrm{i.i.d.}}{\sim} Q^{2,(m)}.
\]
The testing scheme we develop will reject $H_0$ if the distance between the hierarchical estimators exceeds some threshold. In the following sections, we discuss methods for determining this threshold.





















 

